\chapter{Fisika Komputasi}
\label{ch:fiskom}

Kuliah fisika komputasi saya ikuti pada musim dingin 2025/26. Catatan kuliahnya berupa tulisan
tangan, tiga belas pertemuan yang masing-masing membahas satu alat numerik dasar: interpolasi,
transformasi Fourier dan FFT, mencari akar, persamaan diferensial biasa, metode simplektik, masalah
nilai batas, masalah nilai eigen, persamaan diferensial parsial yang bergantung waktu, metode
split-operator, lalu Monte Carlo. Alat-alat ini tersembunyi di dalam hampir setiap simulasi, dan
juga di dalam banyak algoritme machine learning.

\section{Interpolasi dan Fourier}

Interpolasi mencari fungsi mulus yang melewati titik-titik data yang diketahui. Polinomial
interpolasi Lagrange berderajat $n-1$ melewati $n$ titik dengan tepat, tetapi untuk derajat tinggi
pada titik berjarak sama, ia berosilasi liar di dekat tepi selang. Spline kubik memakai polinomial
derajat tiga per potongan yang disambung mulus, dan biasanya jauh lebih jinak. Interpolasi dengan
polinomial trigonometri, yaitu jumlah sinus dan kosinus, cocok untuk data periodik, dan koefisiennya
tidak lain adalah transformasi Fourier diskret.

Menghitung transformasi Fourier diskret dari $N$ titik secara langsung membutuhkan sekitar $N^2$
operasi. FFT \citep{cooley1965} memecah jumlah itu menjadi dua jumlah setengah ukuran, satu untuk indeks
genap dan satu untuk indeks ganjil, lalu mengulang pemecahan itu secara rekursif. Biayanya turun
menjadi sekitar $N\log_2N$. Untuk satu juta titik, selisihnya kira-kira dari $10^{12}$ ke
$2\times10^7$ operasi. FFT menjadi tulang punggung pemrosesan sinyal, pengolahan citra, solver
spektral untuk PDE, dan neural operator berbasis Fourier di \cref{ch:operator}.

\section{Mencari akar}

Banyak masalah berujung pada mencari $x$ yang memenuhi $f(x)=0$. Metode bagi dua mulai dari selang
tempat $f$ berganti tanda, lalu terus membagi selang menjadi dua dan mempertahankan separuh yang
masih mengandung pergantian tanda. Lambat tetapi pasti: setiap langkah memberi tambahan satu bit
ketelitian. Metode Newton--Raphson memakai garis singgung, $x_{n+1}=x_n-f(x_n)/f'(x_n)$, dan
konvergen secara kuadratik di dekat akar, tetapi bisa melompat jauh kalau tebakan awalnya buruk
atau turunannya hampir nol. Dalam banyak dimensi, Newton menjadi penyelesaian sistem linear dengan
Jacobian di setiap langkah, dan itulah yang terjadi di dalam setiap solver implisit.

\section{Persamaan diferensial biasa}

Masalah nilai awal $\dot{\bm y}=\bm f(\bm y,t)$ dengan $\bm y(0)$ diketahui diselesaikan selangkah demi
selangkah. Metode Euler eksplisit, $\bm y_{n+1}=\bm y_n+h\bm f(\bm y_n,t_n)$, hanya akurat orde satu.
Runge--Kutta orde empat mengevaluasi $\bm f$ empat kali per langkah dan galat globalnya sebanding
dengan $h^4$, sehingga memperkecil langkah setengah kali mengurangi galat enam belas kali.

Osilator harmonik memperlihatkan bahwa akurasi saja tidak cukup. Persamaan Newton $\ddot x=-\omega^2x$
punya energi yang kekal, dan di ruang fase posisi dan kecepatan, lintasannya adalah elips tertutup.
Euler eksplisit membuat lintasan berputar keluar, karena energinya bertambah sedikit di setiap
langkah. Runge--Kutta orde empat jauh lebih akurat, tetapi untuk simulasi yang panjang sekali energinya
tetap perlahan bergeser. Metode \istilah{simplektik} seperti velocity Verlet \citep{verlet1967} tidak
menjaga energi persis, tetapi galat energinya berosilasi dalam batas yang tetap dan tidak menumpuk,
karena metode ini menjaga struktur geometris ruang fase, yaitu volume ruang fase, sesuai teorema
Liouville. Untuk dinamika molekul dan mekanika benda langit yang berjalan jutaan langkah, sifat ini
jauh lebih penting daripada orde akurasi.

Masalah nilai batas berbeda. Syaratnya diberikan di dua ujung, misalnya $u(a)=u_a$ dan $u(b)=u_b$.
Metode \istilah{shooting} menebak kemiringan awal, menyelesaikan masalah nilai awal, melihat seberapa
jauh hasilnya meleset di ujung lain, lalu memperbaiki tebakan dengan pencarian akar. Alternatifnya,
turunan diganti beda hingga di grid, sehingga masalahnya menjadi sistem persamaan linear tridiagonal.

\section{Nilai eigen}

Persamaan Schrödinger stasioner dalam satu dimensi adalah masalah nilai eigen: cari energi $E$ dan
fungsi gelombang $\psi$ yang memenuhi $-\tfrac{\hbar^2}{2m}\psi''+V\psi=E\psi$. Dengan beda hingga,
operator turunan kedua menjadi matriks tridiagonal, dan masalahnya menjadi mencari nilai eigen sebuah
matriks. Algoritme QR melakukan ini secara iteratif. Matriks diuraikan menjadi hasil kali matriks
ortogonal dan matriks segitiga atas, dikalikan dalam urutan terbalik, lalu diulang. Untuk matriks
simetris, iterasinya konvergen ke matriks diagonal yang berisi nilai-nilai eigen. Penguraian QR
sendiri dibangun dengan refleksi Householder yang mengenolkan kolom demi kolom.

Nilai eigen muncul di mana-mana di buku ini: arah utama PCA, kestabilan sistem linear, frekuensi
getaran struktur, dan mode-mode aliran pada dekomposisi POD.

\section{PDE yang bergantung waktu}

Persamaan difusi $\partial_tu=D\,\partial_{xx}u$ yang didiskretisasi dengan beda maju dalam waktu dan
beda pusat dalam ruang memberi skema eksplisit
\begin{equation}
u_j^{n+1}=u_j^n+\frac{D\Delta t}{\Delta x^2}\big(u_{j+1}^n-2u_j^n+u_{j-1}^n\big).
\end{equation}
Analisis kestabilan von Neumann memasukkan satu gelombang Fourier $u_j^n=g^ne^{ikj\Delta x}$ ke dalam
skema dan menghitung faktor penguat $g=1-4r\sin^2(k\Delta x/2)$ dengan $r=D\Delta t/\Delta x^2$. Supaya
tidak ada gelombang yang tumbuh, $|g|\le1$ untuk semua $k$, yang berarti $r\le1/2$. Dengan $D=1$ dan
$\Delta x=0{,}01$, langkah waktu tidak boleh melebihi $5\times10^{-5}$. Memperhalus grid dua kali
memaksa langkah waktu empat kali lebih kecil. Skema implisit menghapus batas ini dengan harga
menyelesaikan sistem linear di setiap langkah.

Untuk persamaan Schrödinger yang bergantung waktu, kuliah membahas metode \istilah{split-operator}.
Operator Hamilton terdiri atas energi kinetik, yang sederhana di ruang momentum, dan energi potensial,
yang sederhana di ruang posisi. Satu langkah waktu dipecah menjadi setengah langkah potensial, satu
langkah kinetik setelah FFT ke ruang momentum, lalu setengah langkah potensial lagi setelah FFT balik.
Gagasan memecah operator menjadi bagian-bagian yang masing-masing mudah juga dipakai di banyak solver
fluida, misalnya untuk memisahkan konveksi dari difusi.

\section{Monte Carlo}

Integral berdimensi tinggi tidak bisa dihitung dengan grid, karena jumlah titiknya tumbuh eksponensial.
Integrasi Monte Carlo merata-ratakan fungsi pada titik-titik acak. Galatnya turun seperti $1/\sqrt N$,
lambat, tetapi tidak bergantung pada dimensi sama sekali. Itulah sebabnya Monte Carlo menang untuk
masalah berdimensi tinggi, dan itulah juga hubungannya dengan kutukan dimensi di \cref{ch:data}.

Dalam fisika statistik, kita ingin rata-rata besaran atas distribusi Boltzmann
$p(\bm X)\propto\exp\big(-H(\bm X)/k_BT\big)$, dengan $H$ energi total sistem banyak partikel.
Konstanta normalisasinya, fungsi partisi, mustahil dihitung. Algoritme Metropolis
\citep{metropolis1953} tidak membutuhkannya. Usulkan gerakan kecil acak sebuah partikel. Kalau
energinya turun, terima. Kalau naik sebesar $\Delta H$, terima dengan peluang $\exp(-\Delta H/k_BT)$.
Rantai Markov yang dihasilkan mengunjungi konfigurasi dengan frekuensi sesuai distribusi Boltzmann.
Ini adalah MCMC yang sama dengan di \cref{ch:prob}, di tempat kelahirannya.

\section{Simulasi partikel}

Fisika komputasi juga membawa kita ke simulasi partikel. Dinamika molekul \citep{alder1959} mengintegrasikan
persamaan Newton untuk setiap atom dengan gaya dari potensial antaratom, biasanya dengan velocity
Verlet, dan menghitung sifat makroskopis seperti tekanan dan temperatur sebagai rata-rata
\citep{frenkel2002}. Metode elemen diskret melakukan hal serupa untuk butiran yang bertumbukan
(\cref{ch:numerik}). \istilah{Smoothed particle hydrodynamics} \citep{gingold1977,monaghan1992} memakai
partikel untuk mewakili fluida itu sendiri, dengan besaran di setiap titik diperoleh dengan
merata-ratakan partikel tetangga menggunakan kernel penghalus. Tanpa grid, metode ini nyaman untuk
permukaan bebas dan percikan.

Semua simulasi partikel punya struktur yang sama: setiap partikel berinteraksi dengan tetangga di
dekatnya. Struktur ini persis struktur graf, dan simulator berbasis jaringan graf di \cref{ch:operator}
belajar meniru simulasi partikel seperti ini dengan memperlakukan partikel sebagai simpul dan
interaksi sebagai sisi.

\sumberbab{Bab ini mengikuti tiga belas catatan kuliah Computational Physics I yang ditulis tangan.
Salah satu buku yang disebut di awal kuliah adalah Numerical Recipes \citep{press2007}. Untuk simulasi partikel,
\citet{frenkel2002} adalah pengantar yang baik.}
